\begin{lem} \label{lem:Z2_coefficient}
For any $\phi \colon G \to \Z_2$, there is a split exact sequence of groups
\begin{gather*}
0 \to H^n_G(X; \Z_\phi) \otimes \Z_2 \to H^n_G(X; \Z_2) \to \operatorname{Tor}\big(H^{n+1}_G(X; \Z_\phi), \Z_2\big) \to 0.
\end{gather*}
\end{lem}

\begin{proof}
For any homomorphism $\phi \colon G \to \Z_2$, let $(\Z_2)_\phi$ be the $G$-module such that its underlying group is $\Z_2$ and its $G$-action is given by $\phi \colon G \to \Z_2$. This $G$-module $(\Z_2)_\phi$ agrees with the trivial $G$-module $\Z_2$, even if $\phi$ is non-trivial. Then, looking at the cochain complexes def\/ining the equivariant cohomology, the usual proof of the universal coef\/f\/icient theorem leads to the lemma. Another proof is to use the Thom isomorphism, which unwinds the local coef\/f\/icients: Let $\underline{\R}_{\phi} \to X$ be the $G$-equivariant real line bundle on $X$ whose underlying bundle is $X \times \R$ and the action of $g \in G$ on $(x, r) \in X \times \R$ is $(x, r) \mapsto (gx, \phi(g)r)$. The Thom isomorphism theorem then provides
\begin{gather*}
H^n_G(X; A_\phi) \cong H^{n+1}_G(D, S; A),
\end{gather*}
where $A$ is $\Z_2$ or $\Z$, and $D \subset \underline{\R}_\phi$ and $S \subset \underline{\R}_\phi$ are the unit disk bundle and the unit sphere bundle, respectively. Then the usual universal coef\/f\/icient theorem leads to the present lemma.
\end{proof}

To compute the equivariant cohomology $H^n_G(X; \Z_\phi)$, we usually need the cohomology of the point $H^n_G(\pt; \Z_\phi)$. This cohomology is identif\/ied with the group cohomology $H^n_{\mathrm{group}}(G; \Z_\phi)$ by the degeneration of the Leray--Serre spectral sequence, but its direct computation is not realistic except for the simplest cases (cf.\ Lemma~\ref{lem:group_coh_orientation_coeff}). A useful way to compute it is:

\begin{lem} \label{lem:exact_seq}
For any $\phi \colon G \to \Z_2$, there are natural exact sequences
\begin{gather*}
\cdots \to H^{n-1}_G(\pt; \Z_\phi) \to H^n_G(\pt; \Z) \overset{i^*}{\to} H^n_{\operatorname{Ker}\phi}(\pt; \Z) \to H^n_G(\pt; \Z_\phi) \to \cdots, \\
\cdots \to H^{n-1}_G(\pt; \Z) \to H^n_G(\pt; \Z_\phi) \overset{i^*}{\to} H^n_{\operatorname{Ker}\phi}(\pt; \Z) \to H^n_G(\pt; \Z) \to \cdots,
\end{gather*}
where $i^*$ is induced from the inclusion $i \colon \operatorname{Ker}\phi \to G$.
\end{lem}

\begin{proof}
Let $G$ act on $\R_\phi = \R$ via $\phi \colon G \to \Z_2$. We can regard $\R_\phi$ as a $G$-equivariant real line bundle on $\pt$. We have the Thom isomorphisms
\begin{gather*}
H^n_G(\pt; \Z_\phi) \cong H^{n+1}_G(D, S; \Z), \qquad H^n_G(\pt; \Z) \cong H^{n+1}_G(D, S; \Z_\phi),
\end{gather*}
where $D$ and $S$ are the unit interval in $\R_\phi$ and its boundary, respectively. Note that $D$ is equivariantly contractible. Note also that $S \cong G/\operatorname{Ker}\phi$ as a $G$-space. Thus, we have isomorphisms
\begin{gather*}
H^n_G(D; \Z_\phi) \cong H^n_G(\pt; \Z_\phi), \qquad H^n_G(S; \Z_\phi) \cong H^n_{\operatorname{Ker}\phi}(\pt; \Z).
\end{gather*}
Substituting these isomorphisms and the Thom isomorphisms into the exact sequences for the pair $(D, S)$, we complete the proof.
\end{proof}

By means of the lemma above, we get:

\begin{lem} \label{lem:cohomology_of_point_twisted_case}
Let $P$ be $\Z_{2m}$, $D_{2m-1}$ or $D_{2m}$ with $m \ge 1$. The $P$-equivariant cohomology of the point with coefficients in $\Z_\phi$,
\begin{gather*}
H^n_P(\pt; \Z_\phi) \cong H^n_{\mathrm{group}}(P; \Z_\phi),
\end{gather*}
in low degrees is given as follows:
$$
\begin{array}{|c|c|c|c|c|c|}
\hline
P & \phi &
H^0_P(\pt; \Z_\phi) & H^1_P(\pt; \Z_\phi) &
H^2_P(\pt; \Z_\phi) & H^3_P(\pt; \Z_\phi) \tsep{2pt}\bsep{2pt}\\
\hline
\Z_{2m} & \phi_1 &
0 & \Z_2 & 0 & \Z_2 \\
\hline
D_{2m-1} & \phi_0 &
0 & \Z_2 & \Z_{2m-1} & \Z_2 \\
\hline
D_{2m} & \phi_0 &
0 & \Z_2 & \Z_{2m} & \Z_2^{\oplus 2} \tsep{2pt}\bsep{2pt}\\
\hline
D_{2m} & \phi_1, \phi_2 &
0 & \Z_2 & \Z_2 & \Z_2^{\oplus 2} \tsep{2pt}\bsep{2pt}\\
\hline
\end{array}
$$
\end{lem}

\begin{proof}
In the case of $D_{2m}$ with $m$ even and $\phi \neq \phi_0$, the f\/irst exact sequence in Lemma~\ref{lem:exact_seq} leads to $H^0_{D_{2m}}(\pt; \Z_\phi) = 0$ and $H^1_{D_{2m}}(\pt; \Z_\phi) \cong \Z_2$. This computation also shows that $H^2_{D_{2m}}(\pt; \Z_\phi)$ contains $\Z_2$ as a subgroup. Here, applying the universal coef\/f\/icient theorem to $H^n_P(\pt; \Z)$, we compute the cohomology with coef\/f\/icients in $\Z_2$ to have $H^1_{D_{2m}}(\pt; \Z_2) \cong \Z_2^{\oplus 2}$. If we apply the universal coef\/f\/icient theorem in Lemma \ref{lem:Z2_coefficient} to $H^n_P(\pt; \Z_\phi)$, then
\begin{gather*}
H^1_{D_{2m}}(\pt; \Z_2) \cong \Z_2 \oplus \operatorname{Tor}\big(H^2_{D_{2m}}(\pt; \Z_\phi), \Z_2\big).
\end{gather*}
Thus the consistency of these computations implies $H^2_{D_{2m}}(\pt; \Z_\phi) \cong \Z_2$. Based on this result, the second sequence in Lemma~\ref{lem:exact_seq} suggests that $H^3_{D_{2m}}(\pt; \Z_\phi)$ is either $\Z_2^{\oplus 2}$ or $\Z_4$. If we apply the universal coef\/f\/icient theorem to $H^n_P(\pt; \Z)$, then $H^2_{D_{2m}}(\pt; \Z_2) \cong \Z_2^{\oplus 3}$. If we compute this cohomology applying Lemma~\ref{lem:Z2_coefficient} to $H^n_P(\pt; \Z_\phi)$, then
\begin{gather*}
H^2_{D_{2m}}(\pt; \Z_2) \cong \Z_2 \oplus \operatorname{Tor}\big(H^3_{D_{2m}}(\pt; \Z_\phi), \Z_2\big).
\end{gather*}
Therefore we conclude that $H^3_{D_{2m}}(\pt; \Z_\phi), \Z_2) \cong \Z_2^{\oplus 2}$ by the consistency. In the other cases, a~combined use of the two exact sequences in Lemma~\ref{lem:exact_seq} determines the group $H^n_P(\pt; \Z_\phi)$ without dif\/f\/iculty.
\end{proof}
